% Einheit 1 von 3 – Ansatz und Punktbedingungen (bank/rekonstruktion-von-funktionsgleichungen/e1.jsonl, zusammenbau v0.1)
%% TODO Zeitmarke Einheit 1: Marken-Zeile nicht lesbar
%% TODO Zweigzeile Teil 1: was hier gelernt wird (Satz in Schülersprache aus dem Einheitstitel) – Einheit 1
\einheitenkopf[e1]{Einheit 1 von 3 · Ansatz und Punktbedingungen}
\zweigzeile{Abitur GK}
\setcounter{aufgabe}{8}

%% TODO Titel: Ich-kann-Satz fehlt in der Bank (Platzhalter: Kettenname) – Nr. 9
%% TODO Anweisung: ein Satz über der ersten Teilaufgabe fehlt in der Bank – Nr. 9
\begin{aufgabe}{Ansatz und Punktbedingungen}
%% TODO \rechenplatz{n}: Schrittzahl des Grundfalls fehlt in der Bank (nur bei fester Schreibform, 2.3 b)
\begin{teile}
\teil Gesucht ist eine quadratische Funktion $f(x) = ax^2 + bx + c$; bekannt sind drei Punkte des Graphen. Reichen die Angaben? \\ \kreuz{zu wenig} \\ \kreuz{genau passend} \\ \kreuz{mehr als nötig}
\teil Bestimme die Gleichung der Geraden durch $A(1|5)$ und $B(4|14)$. y = \leerfeld
\teil Die Gerade $h$ verläuft durch $C(0{,}6|1{,}2)$ und $D(3|6)$. Weise nach, dass $h(x) = 2x$ eine Gleichung von $h$ ist. \hfill (Abitur 2022 GK)
\teil Ein Uferweg an einem Stausee wird in einem Koordinatensystem (eine Einheit entspricht 1 km) durch den Graphen einer quadratischen Funktion $g$ beschrieben. Der Graph schneidet die $y$-Achse bei $y = 2$ und verläuft durch $P(1|-1)$ und $Q(3|5)$. Bestimme eine Funktionsgleichung von $g$. \hfill (FHR 2019) \\ g(x) = \leerfeld
\teil Der Graph einer quadratischen Funktion $g$ verläuft symmetrisch zur $y$-Achse; die Punkte $A(1|-3)$ und $B(3|13)$ liegen auf ihm. Bestimme eine Funktionsgleichung von $g$. \hfill (FHR 2023) \\ g(x) = \leerfeld
\teil Die Abbildung zeigt den Graphen einer ganzrationalen Funktion $f$ dritten Grades; er verläuft durch $P(0|3)$. Bestimme einen Funktionsterm von $f$. \hfill (Abitur 2018 LK) \\

\begin{ksys}[xmin=-3,xmax=4,ymin=-3,ymax=5,ablesen] \funktionab{0.5*(\x+2)*(\x-1)*(\x-3)}{f}{-2.3}{3.5} \punkt{0}{3}{P} \end{ksys}
f(x) = \leerfeld
\teil Die Abbildung zeigt den Graphen von $f$ mit $f(x) = a \cdot b^x$ und $a, b > 0$. Lies zwei geeignete Punkte ab und bestimme $a$ und $b$. \hfill (Abitur 2022 GK) \\

\begin{ksys}[xmin=-3,xmax=2,ymin=0,ymax=7,ablesen] \funktionab{2*3^\x}{f}{-3}{1.1} \end{ksys}
a = \leerfeld , b = \leerfeld
\teil Ein neuer Kellergang soll einen parabelförmigen Querschnitt bekommen, 2 m breit und 1,8 m hoch. Im Koordinatensystem liegt der Boden auf der $x$-Achse und der höchste Punkt auf der $y$-Achse; eine Einheit entspricht 1 m. Lies die drei markanten Punkte ab und bestimme eine Gleichung der quadratischen Funktion $g$, deren Graph den Querschnitt beschreibt. \hfill (FHR 2025) \\

\begin{ksys}[xmin=-2,xmax=2,ymin=0,ymax=2,ablesen] \funktionab{-1.8*\x^2+1.8}{g}{-1.0}{1.0} \end{ksys}
Punkte: \leerfeld ; g(x) = \leerfeld
\teil Der Querschnitt einer Futterrinne wird unten durch den Graphen einer Funktion $f$ vierten Grades beschrieben. Der Graph ist achsensymmetrisch zur $y$-Achse, verläuft durch $P(0|-2)$ und hat den Tiefpunkt $T(2|-6)$. Bestimme eine Funktionsgleichung von $f$. \hfill (FHR 2021) \\ f(x) = \leerfeld
\end{teile}
\end{aufgabe}

%% TODO Titel: Ich-kann-Satz fehlt in der Bank (Platzhalter: Kettenname) – Nr. 10
%% TODO Anweisung: ein Satz über der ersten Teilaufgabe fehlt in der Bank – Nr. 10
\begin{aufgabe}{Ansatz und Punktbedingungen – Fehler finden, Begründen}
\begin{teile}
\teil Ich finde den Fehler: Der Graph einer quadratischen Funktion $f$ ist symmetrisch zur $y$-Achse und verläuft durch $A(1|4)$ und $B(2|13)$. Tim setzt $f(x) = ax^2 + bx + c$ an, erhält zwei Gleichungen für drei Unbekannte und schreibt: „Nicht lösbar.“
\teil Begründe, warum beim Ansatz für eine zur $y$-Achse symmetrische Funktion vierten Grades die Glieder mit $x^3$ und $x$ fehlen.
\end{teile}
\end{aufgabe}
